\documentclass[11pt,a4paper]{article}

\usepackage[utf8]{inputenc}
\usepackage[T1]{fontenc}
\usepackage[english]{babel}
\usepackage[margin=2.5cm]{geometry}
\usepackage{amsmath}
\usepackage{amssymb}
\usepackage{amsthm}
\usepackage{mathtools}
\usepackage{microtype}
\usepackage{hyperref}
\usepackage{cleveref}
\usepackage{fancyhdr}
\usepackage{titlesec}
\usepackage{abstract}

% Hyperref setup
\hypersetup{
    colorlinks=true,
    linkcolor=blue,
    filecolor=magenta,      
    urlcolor=cyan,
    citecolor=red,
    pdftitle={The Neural Tangent Kernel and its First-Order Corrections},
    pdfauthor={Research Report},
    pdfsubject={Neural Networks, NTK Theory},
    pdfkeywords={Neural Tangent Kernel, First-order corrections, ReLU networks}
}

% Page style
\pagestyle{fancy}
\fancyhf{}
\fancyhead[L]{\leftmark}
\fancyhead[R]{\thepage}
\fancyfoot[C]{\footnotesize Neural Tangent Kernel Theory}

% Theorem environments
\newtheorem{theorem}{Theorem}[section]
\newtheorem{proposition}[theorem]{Proposition}
\newtheorem{lemma}[theorem]{Lemma}
\newtheorem{corollary}[theorem]{Corollary}
\newtheorem{remark}[theorem]{Remark}
\newtheorem{definition}[theorem]{Definition}

\theoremstyle{definition}
\newtheorem{example}[theorem]{Example}

% Title formatting
\titleformat{\section}
  {\normalfont\Large\bfseries}{\thesection}{1em}{}
\titleformat{\subsection}
  {\normalfont\large\bfseries}{\thesubsection}{1em}{}

% Abstract formatting
\renewcommand{\abstractname}{Abstract}
\renewcommand{\absnamepos}{center}

\title{\LARGE\textbf{The Neural Tangent Kernel and its First-Order Corrections}\\
\vspace{0.5cm}
\large A Mathematical Analysis of Finite-Width Networks}

\author{
\textbf{Research Report}\\
\vspace{0.3cm}
\textit{Neural Network Theory Group}\\
\vspace{0.2cm}
\today
}

\date{}

\begin{document}

\maketitle

\begin{abstract}
This paper presents a comprehensive mathematical analysis of the Neural Tangent Kernel (NTK) and its first-order corrections for finite-width neural networks. We provide a detailed derivation showing how the corrections $O_3$ and $O_4$ emerge from the gradient of the NTK formula with respect to network parameters. Our main contribution is a rigorous proof demonstrating that for ReLU activation functions, the backpropagation terms containing second derivatives vanish almost everywhere, leading to a natural truncation of the correction series. We establish a spectral decomposition of the NTK eigenvalues, showing that the minimal eigenvalue admits the form $\lambda_{\min}(K_\theta) \geq al + \frac{b}{N} + o(\frac{1}{N})$, where the linear term $al$ represents the infinite-width limit and $\frac{b}{N}$ captures the finite-width corrections. This analysis provides crucial insights into the convergence properties of gradient descent in finite-width neural networks.
\end{abstract}

\vspace{1cm}

\tableofcontents

\newpage


\section{Other proof}

In the framework of the Neural Tangent Hierarchy (NTH), the kernel $K^{(3)}$ arises as the operator governing the time evolution of the Neural Tangent Kernel $K^{(2)}$. The dynamics of $K^{(2)}$ under gradient flow are given by:
\begin{equation}
\frac{d}{dt} K^{(2)}_t(x_\alpha, x_\beta) = -\frac{1}{n} \sum_{\gamma=1}^n K^{(3)}_t(x_\alpha, x_\beta, x_\gamma) (f_t(x_\gamma) - y_\gamma)
\end{equation}
where $K^{(3)}_t(x_\alpha, x_\beta, x_\gamma)$ is defined as the contraction of the gradient of $K^{(2)}_t(x_\alpha, x_\beta)$ with respect to the network parameters $\theta$ with the gradient of the network output $f_t(x_\gamma)$:
\begin{equation}
K^{(3)}_t(x_\alpha, x_\beta, x_\gamma) := \langle \nabla_\theta K^{(2)}_t(x_\alpha, x_\beta), \nabla_\theta f_t(x_\gamma) \rangle.
\end{equation}

Let's derive the explicit formula for $K^{(3)}_t$. We start from the expression for $K^{(2)}_t(x_\alpha, x_\beta)$:
\begin{equation}
K^{(2)}_t(x_\alpha, x_\beta) = \langle x^{(H)}_\alpha, x^{(H)}_\beta \rangle + \sum_{\ell=1}^{H} \langle G^{(\ell)}_\alpha, G^{(\ell)}_\beta \rangle \langle x^{(\ell-1)}_\alpha, x^{(\ell-1)}_\beta \rangle
\end{equation}
where we use the notation:
\begin{itemize}
    \item $x^{(p)}_\mu = \frac{1}{\sqrt{m}}\sigma(W^{(p)} x^{(p-1)}_\mu)$ for $p \ge 1$, and $x^{(0)}_\mu = x_\mu$.
    \item $G^{(\ell)}_\mu = \frac{1}{\sqrt{m}}\sigma'_\ell(x_\mu) (W^{(\ell+1)})^T \cdots \frac{1}{\sqrt{m}}\sigma'_H(x_\mu) a_t$, with $\sigma'_\ell(x_\mu) = \text{diag}(\sigma'(W^{(\ell)}x^{(\ell-1)}_\mu))$.
\end{itemize}

Let $\delta_\gamma(\cdot) := \langle \nabla_\theta (\cdot), \nabla_\theta f_t(x_\gamma) \rangle$ denote the directional derivative along $\nabla_\theta f_t(x_\gamma)$. Then $K^{(3)}_{\alpha\beta\gamma} = \delta_\gamma(K^{(2)}_{\alpha\beta})$. Applying the product rule, we get:
\begin{align}
K^{(3)}_{\alpha\beta\gamma} = \delta_\gamma \langle x^{(H)}_\alpha, x^{(H)}_\beta \rangle + \sum_{\ell=1}^{H} \delta_\gamma \left( \langle G^{(\ell)}_\alpha, G^{(\ell)}_\beta \rangle \langle x^{(\ell-1)}_\alpha, x^{(\ell-1)}_\beta \rangle \right)
\end{align}
Each term expands further, following $\delta_\gamma \langle A, B \rangle = \langle \delta_\gamma A, B \rangle + \langle A, \delta_\gamma B \rangle$. This results in a sum over all components of $K^{(2)}_{\alpha\beta}$, where each component is acted upon by $\delta_\gamma$. The core of the derivation lies in computing $\delta_\gamma$ for the building blocks $x^{(p)}_\mu$ and $G^{(p)}_\mu$.

The recursive formula for the derivative of the activations is:
\begin{equation}
\delta_\gamma x^{(p)}_\mu = \frac{1}{\sqrt{m}} \sigma'_{p}(x_\mu) \left( W^{(p)} (\delta_\gamma x^{(p-1)}_\mu) + \langle x^{(p-1)}_\gamma, x^{(p-1)}_\mu \rangle G^{(p)}_\gamma \right)
\end{equation}
with the base case $\delta_\gamma x^{(0)}_\mu = \delta_\gamma x_\mu = 0$.

The derivative of the backpropagated vectors $G^{(\ell)}_\mu$ is given by a sum of terms, where $\delta_\gamma$ acts on one factor at a time:
\begin{align}
\delta_\gamma G^{(\ell)}_\mu = \sum_{p=\ell}^{H} & \left( \frac{1}{\sqrt{m}}\sigma'_\ell(x_\mu) \cdots \delta_\gamma(\frac{1}{\sqrt{m}}(W^{(p+1)})^T) \cdots \sigma'_H(x_\mu) a_t \right) \\
+ \sum_{p=\ell}^{H} & \left( \frac{1}{\sqrt{m}}\sigma'_\ell(x_\mu) \cdots \delta_\gamma(\sigma'_{p}(x_\mu)) \cdots \sigma'_H(x_\mu) a_t \right) \\
+ & \left( \frac{1}{\sqrt{m}}\sigma'_\ell(x_\mu) \cdots \sigma'_H(x_\mu) (\delta_\gamma a_t) \right)
\end{align}
where
\begin{itemize}
    \item $\delta_\gamma a_t = x^{(H)}_\gamma$
    \item The term $\delta_\gamma((W^{(p+1)})^T)$ corresponds to a rank-1 update which, when contracted in context, yields terms involving $G^{(p+1)}_\gamma$ and inner products of activations.
    \item $\delta_\gamma(\sigma'_{p}(x_\mu)) = \sigma''_{p}(x_\mu) \text{diag}(\delta_\gamma(W^{(p)}x^{(\ell-1)}_\mu/\sqrt{m}))$. This term vanishes for ReLU activation functions as $\sigma''=0$ almost everywhere.
\end{itemize}

Combining these rules gives the complete formula for $K^{(3)}_t(x_\alpha, x_\beta, x_\gamma)$. It is a sum of many terms, each corresponding to differentiating one part of $K^{(2)}_t(x_\alpha, x_\beta)$. For clarity, we write it as a sum of contributions:
\begin{equation}
K^{(3)}_{\alpha\beta\gamma} = K^{(3, \text{out})}_{\alpha\beta\gamma} + \sum_{\ell=1}^H \left( K^{(3, G)}_{\alpha\beta\gamma, \ell} + K^{(3, x)}_{\alpha\beta\gamma, \ell} \right)
\end{equation}
where
\begin{align}
K^{(3, \text{out})}_{\alpha\beta\gamma} &= \langle \delta_\gamma x^{(H)}_\alpha, x^{(H)}_\beta \rangle + \langle x^{(H)}_\alpha, \delta_\gamma x^{(H)}_\beta \rangle \\
K^{(3, G)}_{\alpha\beta\gamma, \ell} &= \left( \langle \delta_\gamma G^{(\ell)}_\alpha, G^{(\ell)}_\beta \rangle + \langle G^{(\ell)}_\alpha, \delta_\gamma G^{(\ell)}_\beta \rangle \right) \langle x^{(\ell-1)}_\alpha, x^{(\ell-1)}_\beta \rangle \\
K^{(3, x)}_{\alpha\beta\gamma, \ell} &= \langle G^{(\ell)}_\alpha, G^{(\ell)}_\beta \rangle \left( \langle \delta_\gamma x^{(\ell-1)}_\alpha, x^{(\ell-1)}_\beta \rangle + \langle x^{(\ell-1)}_\alpha, \delta_\gamma x^{(\ell-1)}_\beta \rangle \right)
\end{align}
The terms $\delta_\gamma x^{(p)}_\mu$ and $\delta_\gamma G^{(p)}_\mu$ are expanded using the recursive rules above, yielding a full, albeit lengthy, expression. This hierarchical structure is characteristic of the NTH framework.


\section{Extended Formula for $K^{(3)}$}

To obtain a complete expression for $K^{(3)}_{\alpha\beta\gamma}$, we expand the terms $\delta_\gamma x^{(p)}_\mu$ and $\delta_\gamma G^{(\ell)}_\mu$ using their recursive definitions. We assume a ReLU activation function, so $\sigma''=0$ almost everywhere.

\subsection{Development of Recursive Terms}

The term $\delta_\gamma x^{(p)}_\mu$, which represents the variation of activation $x^{(p)}_\mu$ along the gradient of output $f_t(x_\gamma)$, can be expanded as a sum over previous layers:
\begin{equation}
\delta_\gamma x^{(p)}_\mu = \sum_{j=1}^{p} \left( \prod_{k=j+1}^{p} \frac{\sigma'_{k}(x_\mu) W^{(k)}}{\sqrt{m}} \right) \frac{\sigma'_{j}(x_\mu)}{\sqrt{m}} \langle x^{(j-1)}_\gamma, x^{(j-1)}_\mu \rangle G^{(j)}_\gamma
\end{equation}
where the product $\prod_{k=j+1}^{p}$ is an ordered product of Jacobian matrices.

Similarly, the term $\delta_\gamma G^{(\ell)}_\mu$, the variation of the backpropagated gradient vector, decomposes into a sum over subsequent layers:
\begin{align}
\delta_\gamma G^{(\ell)}_\mu = & \sum_{p=\ell}^{H-1} \left( \prod_{k=\ell}^{p} \frac{(W^{(k+1)})^T \sigma'_{k+1}(x_\mu)}{\sqrt{m}} \right) \frac{1}{m} G^{(p+1)}_\gamma \langle x^{(p)}_\gamma, x^{(p)}_\mu \rangle \\
& + \left( \prod_{k=\ell}^{H-1} \frac{(W^{(k+1)})^T \sigma'_{k+1}(x_\mu)}{\sqrt{m}} \right) \frac{x^{(H)}_\gamma}{\sqrt{m}}
\end{align}
The first term captures the effect of varying weights $W^{(\ell+1)}, \dots, W^{(H)}$, and the second term that of varying the output vector $a_t$.

\subsection{Complete Expression}

Substituting these expressions into equations (18), (19) and (20) yields the complete formula for $K^{(3)}_{\alpha\beta\gamma}$. It is a sum of $2H^2 + 2H$ terms.
\begin{align}
K^{(3)}_{\alpha\beta\gamma} = & \sum_{j=1}^{H} \left( \langle \mathcal{J}^{(H \to j)}_{\alpha} \mathcal{T}^{(j)}_{\gamma\alpha}, x^{(H)}_\beta \rangle + \langle x^{(H)}_\alpha, \mathcal{J}^{(H \to j)}_{\beta} \mathcal{T}^{(j)}_{\gamma\beta} \rangle \right) \nonumber \\
& + \sum_{\ell=1}^{H} \langle x^{(\ell-1)}_\alpha, x^{(\ell-1)}_\beta \rangle \sum_{p=\ell}^{H-1} \left( \langle \mathcal{B}^{(\ell \to p)}_{\alpha} \mathcal{U}^{(p)}_{\gamma\alpha}, G^{(\ell)}_\beta \rangle + \langle G^{(\ell)}_\alpha, \mathcal{B}^{(\ell \to p)}_{\beta} \mathcal{U}^{(p)}_{\gamma\beta} \rangle \right) \nonumber \\
& + \sum_{\ell=1}^{H} \langle x^{(\ell-1)}_\alpha, x^{(\ell-1)}_\beta \rangle \left( \langle \mathcal{B}^{(\ell \to H-1)}_{\alpha} \frac{x^{(H)}_\gamma}{\sqrt{m}}, G^{(\ell)}_\beta \rangle + \langle G^{(\ell)}_\alpha, \mathcal{B}^{(\ell \to H-1)}_{\beta} \frac{x^{(H)}_\gamma}{\sqrt{m}} \rangle \right) \nonumber \\
& + \sum_{\ell=1}^{H} \langle G^{(\ell)}_\alpha, G^{(\ell)}_\beta \rangle \sum_{j=1}^{\ell-1} \left( \langle \mathcal{J}^{(\ell-1 \to j)}_{\alpha} \mathcal{T}^{(j)}_{\gamma\alpha}, x^{(\ell-1)}_\beta \rangle + \langle x^{(\ell-1)}_\alpha, \mathcal{J}^{(\ell-1 \to j)}_{\beta} \mathcal{T}^{(j)}_{\gamma\beta} \rangle \right)
\end{align}
where we have defined for readability:
\begin{itemize}
    \item Forward propagators (Jacobian): $\mathcal{J}^{(p \to j)}_{\mu} = \prod_{k=j+1}^{p} \frac{\sigma'_{k}(x_\mu) W^{(k)}}{\sqrt{m}}$
    \item Forward source terms: $\mathcal{T}^{(j)}_{\gamma\mu} = \frac{\sigma'_{j}(x_\mu)}{\sqrt{m}} \langle x^{(j-1)}_\gamma, x^{(j-1)}_\mu \rangle G^{(j)}_\gamma$
    \item Backward propagators: $\mathcal{B}^{(\ell \to p)}_{\mu} = \prod_{k=\ell}^{p} \frac{(W^{(k+1)})^T \sigma'_{k+1}(x_\mu)}{\sqrt{m}}$
    \item Backward source terms: $\mathcal{U}^{(p)}_{\gamma\mu} = \frac{1}{m} G^{(p+1)}_\gamma \langle x^{(p)}_\gamma, x^{(p)}_\mu \rangle$
\end{itemize}

Using the readability notations defined above, we can write the components of $K^{(3)}_{\alpha\beta\gamma}$ in a more compact form:

\begin{align}
K^{(3, \text{out})}_{\alpha\beta\gamma} &= \sum_{j=1}^{H} \left( \langle \mathcal{J}^{(H \to j)}_{\alpha} \mathcal{T}^{(j)}_{\gamma\alpha}, x^{(H)}_\beta \rangle + \langle x^{(H)}_\alpha, \mathcal{J}^{(H \to j)}_{\beta} \mathcal{T}^{(j)}_{\gamma\beta} \rangle \right) \\
K^{(3, G)}_{\alpha\beta\gamma, \ell} &= \langle x^{(\ell-1)}_\alpha, x^{(\ell-1)}_\beta \rangle \sum_{p=\ell}^{H-1} \left( \langle \mathcal{B}^{(\ell \to p)}_{\alpha} \mathcal{U}^{(p)}_{\gamma\alpha}, G^{(\ell)}_\beta \rangle + \langle G^{(\ell)}_\alpha, \mathcal{B}^{(\ell \to p)}_{\beta} \mathcal{U}^{(p)}_{\gamma\beta} \rangle \right) \nonumber \\
&\quad + \langle x^{(\ell-1)}_\alpha, x^{(\ell-1)}_\beta \rangle \left( \langle \mathcal{B}^{(\ell \to H-1)}_{\alpha} \frac{x^{(H)}_\gamma}{\sqrt{m}}, G^{(\ell)}_\beta \rangle + \langle G^{(\ell)}_\alpha, \mathcal{B}^{(\ell \to H-1)}_{\beta} \frac{x^{(H)}_\gamma}{\sqrt{m}} \rangle \right) \\
K^{(3, x)}_{\alpha\beta\gamma, \ell} &= \langle G^{(\ell)}_\alpha, G^{(\ell)}_\beta \rangle \sum_{j=1}^{\ell-1} \left( \langle \mathcal{J}^{(\ell-1 \to j)}_{\alpha} \mathcal{T}^{(j)}_{\gamma\alpha}, x^{(\ell-1)}_\beta \rangle + \langle x^{(\ell-1)}_\alpha, \mathcal{J}^{(\ell-1 \to j)}_{\beta} \mathcal{T}^{(j)}_{\gamma\beta} \rangle \right)
\end{align}
and 
\begin{equation}
    K^{(3)}_{\alpha\beta\gamma} = K^{(3, \text{out})}_{\alpha\beta\gamma} + \sum_{\ell=1}^H \left( K^{(3, G)}_{\alpha\beta\gamma, \ell} + K^{(3, x)}_{\alpha\beta\gamma, \ell} \right)
    \end{equation}

This representation makes the hierarchical structure more apparent, with clear separation between forward and backward propagation terms, and highlights the symmetry between the $\alpha$ and $\beta$ indices.


This formula, while complex, represents the complete hierarchical structure of the first-order correction to the NTK dynamics.


\section{Fully Expanded Non-Recursive Formula for $K^{(3)}$}

Here, we present the complete, non-recursive expression for $K^{(3)}_{\alpha\beta\gamma}$. This formula is obtained by substituting the expanded expressions for the directional derivatives $\delta_\gamma x^{(p)}_\mu$ and $\delta_\gamma G^{(\ell)}_\mu$ into the decomposed formula for $K^{(3)}$. This makes the dependency on the network parameters (weights $W^{(k)}$ and output vectors $a_t$) and activations at each layer fully explicit. We continue to assume a ReLU activation function, which sets second derivatives of $\sigma$ to zero almost everywhere.

The final expression is a sum of four main components, corresponding to the differentiation of the output layer activations, the backward vectors, and the hidden layer activations:

\begin{align}
K^{(3)}_{\alpha\beta\gamma} = & \nonumber \\
% K^(3,out) term
& \sum_{j=1}^{H} \Biggl( \left\langle \left( \prod_{k=j+1}^{H} \frac{\sigma'_{k}(x_\alpha) W^{(k)}}{\sqrt{m}} \right) \frac{\sigma'_{j}(x_\alpha)}{\sqrt{m}} \langle x^{(j-1)}_\gamma, x^{(j-1)}_\alpha \rangle G^{(j)}_\gamma, x^{(H)}_\beta \right\rangle \nonumber \\
& \quad + \left\langle x^{(H)}_\alpha, \left( \prod_{k=j+1}^{H} \frac{\sigma'_{k}(x_\beta) W^{(k)}}{\sqrt{m}} \right) \frac{\sigma'_{j}(x_\beta)}{\sqrt{m}} \langle x^{(j-1)}_\gamma, x^{(j-1)}_\beta \rangle G^{(j)}_\gamma \right\rangle \Biggr) \nonumber \\
% K^(3,G) term
& + \sum_{\ell=1}^{H} \langle x^{(\ell-1)}_\alpha, x^{(\ell-1)}_\beta \rangle \Biggl[ \nonumber \\
& \quad \sum_{p=\ell}^{H-1} \Biggl( \left\langle \left( \prod_{k=\ell}^{p} \frac{(W^{(k+1)})^T \sigma'_{k+1}(x_\alpha)}{\sqrt{m}} \right) \frac{G^{(p+1)}_\gamma \langle x^{(p)}_\gamma, x^{(p)}_\alpha \rangle}{m}, G^{(\ell)}_\beta \right\rangle \nonumber \\
& \qquad + \left\langle G^{(\ell)}_\alpha, \left( \prod_{k=\ell}^{p} \frac{(W^{(k+1)})^T \sigma'_{k+1}(x_\beta)}{\sqrt{m}} \right) \frac{G^{(p+1)}_\gamma \langle x^{(p)}_\gamma, x^{(p)}_\beta \rangle}{m} \right\rangle \Biggr) \nonumber \\
& \quad + \left\langle \left( \prod_{k=\ell}^{H-1} \frac{(W^{(k+1)})^T \sigma'_{k+1}(x_\alpha)}{\sqrt{m}} \right) \frac{x^{(H)}_\gamma}{\sqrt{m}}, G^{(\ell)}_\beta \right\rangle \nonumber \\
& \quad + \left\langle G^{(\ell)}_\alpha, \left( \prod_{k=\ell}^{H-1} \frac{(W^{(k+1)})^T \sigma'_{k+1}(x_\beta)}{\sqrt{m}} \right) \frac{x^{(H)}_\gamma}{\sqrt{m}} \right\rangle \Biggr] \nonumber \\
% K^(3,x) term
& + \sum_{\ell=1}^{H} \langle G^{(\ell)}_\alpha, G^{(\ell)}_\beta \rangle \sum_{j=1}^{\ell-1} \Biggl( \left\langle \left( \prod_{k=j+1}^{\ell-1} \frac{\sigma'_{k}(x_\alpha) W^{(k)}}{\sqrt{m}} \right) \frac{\sigma'_{j}(x_\alpha)}{\sqrt{m}} \langle x^{(j-1)}_\gamma, x^{(j-1)}_\alpha \rangle G^{(j)}_\gamma, x^{(\ell-1)}_\beta \right\rangle \nonumber \\
& \quad + \left\langle x^{(\ell-1)}_\alpha, \left( \prod_{k=j+1}^{\ell-1} \frac{\sigma'_{k}(x_\beta) W^{(k)}}{\sqrt{m}} \right) \frac{\sigma'_{j}(x_\beta)}{\sqrt{m}} \langle x^{(j-1)}_\gamma, x^{(j-1)}_\beta \rangle G^{(j)}_\gamma \right\rangle \Biggr)
\end{align}
In this expression, the products $\prod$ represent ordered matrix products. The empty product $\prod_{k=p+1}^p$ is defined as the identity matrix. This detailed formula lays bare the intricate dependencies of the NTK dynamics on the network's architecture and parameters at finite width.


\newpage

\section{Framework for the $K^{(4)}$ Kernel Derivation}

This section lays the groundwork for the derivation of the fourth-order kernel, $K^{(4)}$. Following the same logic as the Neural Tangent Hierarchy, $K^{(4)}$ arises from the dynamics of $K^{(3)}$. Our goal here is not to perform the full derivation, which is algebraically intensive, but to outline the structure of the proof, define the necessary objects, and describe the computational steps involved.

\subsection{Definition of $K^{(4)}$ and the Directional Derivative}

We define the $K^{(4)}$ kernel as the change of the $K^{(3)}$ kernel along the gradient of the network's output with respect to a new data point $x_\mu$.
\begin{definition}[The $K^{(4)}$ Kernel]
The kernel $K^{(4)}_{\alpha\beta\gamma\mu}$ is defined as the inner product of the gradient of $K^{(3)}_{\alpha\beta\gamma}$ with respect to the network parameters $\theta$ and the gradient of the network output $f_t(x_\mu)$:
\begin{equation}
K^{(4)}_{\alpha\beta\gamma\mu} := \langle \nabla_\theta K^{(3)}_{\alpha\beta\gamma}, \nabla_\theta f_t(x_\mu) \rangle
\end{equation}
\end{definition}
To streamline the derivation, we introduce a new directional derivative operator, $\delta_\mu$, acting on any function $F(\theta)$ of the parameters:
\begin{equation}
\delta_\mu(F) := \langle \nabla_\theta F, \nabla_\theta f_t(x_\mu) \rangle
\end{equation}
With this notation, the fourth-order kernel is simply $K^{(4)}_{\alpha\beta\gamma\mu} = \delta_\mu(K^{(3)}_{\alpha\beta\gamma})$.

\subsection{Derivation Strategy}

The derivation of $K^{(4)}$ proceeds by applying the operator $\delta_\mu$ to the complete, non-recursive formula for $K^{(3)}_{\alpha\beta\gamma}$ presented in the previous section. Since $K^{(3)}$ is a sum of terms, we can apply the derivative term by term. The expression for $K^{(3)}$ is composed of fundamental building blocks:
\begin{itemize}
    \item Activations: $x^{(p)}_\nu$
    \item Backpropagated gradient vectors: $G^{(\ell)}_\nu$
    \item Inner products of these vectors.
    \item Forward and backward propagators, $\mathcal{J}$ and $\mathcal{B}$.
\end{itemize}
The core of the computation will involve applying $\delta_\mu$ to these blocks and then using the product rule for inner products and matrix-vector products. For example, for a generic term $\langle A, B \rangle$ in the expression of $K^{(3)}$, its derivative will be $\delta_\mu \langle A, B \rangle = \langle \delta_\mu A, B \rangle + \langle A, \delta_\mu B \rangle$.

\subsection{An Alternative Perspective: A Rewriting System}

An alternative way to conceptualize the computation of the directional derivative is through a "rewriting system". Instead of applying differential calculus rules, we can view the operation $\delta_\gamma$ as a set of symbolic replacement rules applied to every component of the kernel expression. To obtain $K^{(3)}_{\alpha\beta\gamma} = \delta_\gamma(K^{(2)}_{\alpha\beta})$, one takes the formula for $K^{(2)}_{\alpha\beta}$ and sums all possible terms generated by applying one of the following replacements to each of its constituent parts (for $\nu = \alpha, \beta$):

\begin{itemize}
    \item \textbf{Output Layer:} The output vector $a_t$ is replaced by the final layer activation corresponding to the derivative index $\gamma$.
    \begin{equation}
    a_t \rightarrow x^{(H)}_\gamma
    \end{equation}
    \item \textbf{Weight Matrices:} A weight matrix $W^{(k)}$ is replaced by a rank-one matrix representing its gradient contribution.
    \begin{equation}
    W^{(k)} \rightarrow G^{(k)}_\gamma (x^{(k-1)}_\gamma)^T
    \end{equation}
    A similar rule applies to its transpose $(W^{(k)})^T$.
    \item \textbf{Activations:} The activation $x^{(\ell)}_\nu$ is replaced by its variation, which we previously derived.
    \begin{equation}
    x^{(\ell)}_\nu \rightarrow \delta_\gamma x^{(\ell)}_\nu = \sum_{j=1}^{\ell} \left( \prod_{k=j+1}^{\ell} \frac{\sigma'_{k}(x_\nu) W^{(k)}}{\sqrt{m}} \right) \frac{\sigma'_{j}(x_\nu)}{\sqrt{m}} \langle x^{(j-1)}_\gamma, x^{(j-1)}_\nu \rangle G^{(j)}_\gamma
    \end{equation}
    \item \textbf{Activation Derivatives:} The derivative of the activation function $\sigma'_\ell(x_\nu)$ is replaced by a term involving the second derivative $\sigma''_\ell(x_\nu)$.
    \begin{equation}
    \sigma'_\ell(x_\nu) \rightarrow \sigma''_\ell(x_\nu) \text{diag}(\delta_\gamma(W^{(\ell)}x^{(\ell-1)}_\nu/\sqrt{m}))
    \end{equation}
\end{itemize}

\begin{remark}[Simplification for ReLU Networks]
A key simplification arises for ReLU (and other piecewise linear) activation functions. For ReLU, the second derivative $\sigma''$ is zero almost everywhere. Consequently, the replacement rule for $\sigma'_\ell(x_\nu)$ contributes nothing to the final expression. This is precisely why the terms containing second derivatives vanished in our earlier derivation of $K^{(3)}$.
\end{remark}

This rewriting system offers a powerful combinatorial approach. To compute $K^{(4)}_{\alpha\beta\gamma\mu}$, we would simply apply a second set of replacement rules, this time indexed by $\mu$, to the full expression for $K^{(3)}_{\alpha\beta\gamma}$. This provides a second, complementary technique to the direct differentiation method, which can serve as a valuable cross-check for the results.

\subsection{Core Second-Order Derivative Computations}

The main challenge is to compute the action of $\delta_\mu$ on the terms that are already a result of the $\delta_\gamma$ differentiation. This will lead to second-order derivative objects. The essential computations to be performed are:
\begin{enumerate}
    \item \textbf{Derivatives of activations:} We need to compute $\delta_\mu(x^{(p)}_\nu)$. The formula is identical to that of $\delta_\gamma(x^{(p)}_\nu)$, with $\gamma$ replaced by $\mu$.
    \item \textbf{Derivatives of backpropagated vectors:} Similarly, we need $\delta_\mu(G^{(\ell)}_\nu)}$, which follows the same structure as $\delta_\gamma(G^{(\ell)}_\nu)$.
    \item \textbf{Derivatives of propagators:} The propagators $\mathcal{J}^{(p \to j)}_{\nu}$ and $\mathcal{B}^{(\ell \to p)}_{\nu}$ are products of weight matrices and activation derivatives. Their differentiation with $\delta_\mu$ will produce a sum of terms, according to the product rule. For ReLU networks, terms involving $\delta_\mu(\sigma'_k)$ will vanish as they depend on $\sigma''=0$. The main contribution will come from differentiating the weight matrices: $\delta_\mu(W^{(k)})$.
    \item \textbf{Derivatives of source terms:} We must compute derivatives of the source terms $\mathcal{T}^{(j)}_{\gamma\nu}$ and $\mathcal{U}^{(p)}_{\gamma\nu}$. This is where the interaction between the derivatives $\delta_\mu$ and $\delta_\gamma$ will appear. For example, computing $\delta_\mu(\mathcal{T}^{(j)}_{\gamma\nu})$ will require evaluating $\delta_\mu(G^{(j)}_\gamma)$ and $\delta_\mu(\langle x^{(j-1)}_\gamma, x^{(j-1)}_\nu \rangle)$, the latter of which expands into terms containing $\delta_\mu x^{(j-1)}_\gamma$.
\end{enumerate}

\subsection{Notation for Higher-Order Terms}

To manage the complexity of the final expression for $K^{(4)}$, a systematic notation for the second-order terms will be required. We could introduce notations such as:
\begin{itemize}
    \item Second-order forward source terms: $\mathcal{T}^{(j)}_{\gamma\mu\nu} := \delta_\mu(\mathcal{T}^{(j)}_{\gamma\nu})$.
    \item Second-order backward source terms: $\mathcal{U}^{(p)}_{\gamma\mu\nu} := \delta_\mu(\mathcal{U}^{(p)}_{\gamma\nu})$.
\end{itemize}
These objects would encapsulate the result of applying one directional derivative after another. Their full expression would be derived by applying the rules described in the previous subsection. The final formula for $K^{(4)}_{\alpha\beta\gamma\mu}$ would then be expressed as a sum of inner products involving these new higher-order objects, preserving the hierarchical structure observed for $K^{(3)}$.


\newpage

\section{Full Expression for $K^{(4)}$ via rewriting system}

Here, we derive the complete expression for $K^{(4)}_{\alpha\beta\gamma\mu}$ by applying the rewriting system rules, corresponding to the directional derivative $\delta_\mu$, to the full expression of $K^{(3)}_{\alpha\beta\gamma}$.

\subsection{The Rewriting Rules for $\delta_\mu$}

We recall the formula for $K^{(3)}$:
\begin{align}
    K^{(3)}_{\alpha\beta\gamma} = & \sum_{j=1}^{H} \left( \langle \mathcal{J}^{(H \to j)}_{\alpha} \mathcal{T}^{(j)}_{\gamma\alpha}, x^{(H)}_\beta \rangle + \langle x^{(H)}_\alpha, \mathcal{J}^{(H \to j)}_{\beta} \mathcal{T}^{(j)}_{\gamma\beta} \rangle \right) \nonumber \\
    & + \sum_{\ell=1}^{H} \langle x^{(\ell-1)}_\alpha, x^{(\ell-1)}_\beta \rangle \sum_{p=\ell}^{H-1} \left( \langle \mathcal{B}^{(\ell \to p)}_{\alpha} \mathcal{U}^{(p)}_{\gamma\alpha}, G^{(\ell)}_\beta \rangle + \langle G^{(\ell)}_\alpha, \mathcal{B}^{(\ell \to p)}_{\beta} \mathcal{U}^{(p)}_{\gamma\beta} \rangle \right) \nonumber \\
    & + \sum_{\ell=1}^{H} \langle x^{(\ell-1)}_\alpha, x^{(\ell-1)}_\beta \rangle \left( \langle \mathcal{B}^{(\ell \to H-1)}_{\alpha} \frac{x^{(H)}_\gamma}{\sqrt{m}}, G^{(\ell)}_\beta \rangle + \langle G^{(\ell)}_\alpha, \mathcal{B}^{(\ell \to H-1)}_{\beta} \frac{x^{(H)}_\gamma}{\sqrt{m}} \rangle \right) \nonumber \\
    & + \sum_{\ell=1}^{H} \langle G^{(\ell)}_\alpha, G^{(\ell)}_\beta \rangle \sum_{j=1}^{\ell-1} \left( \langle \mathcal{J}^{(\ell-1 \to j)}_{\alpha} \mathcal{T}^{(j)}_{\gamma\alpha}, x^{(\ell-1)}_\beta \rangle + \langle x^{(\ell-1)}_\alpha, \mathcal{J}^{(\ell-1 \to j)}_{\beta} \mathcal{T}^{(j)}_{\gamma\beta} \rangle \right)
    \end{align}
    where we have defined for readability:
    \begin{itemize}
        \item Forward propagators (Jacobian): $\mathcal{J}^{(p \to j)}_{\mu} = \prod_{k=j+1}^{p} \frac{\sigma'_{k}(x_\mu) W^{(k)}}{\sqrt{m}}$
        \item Forward source terms: $\mathcal{T}^{(j)}_{\gamma\mu} = \frac{\sigma'_{j}(x_\mu)}{\sqrt{m}} \langle x^{(j-1)}_\gamma, x^{(j-1)}_\mu \rangle G^{(j)}_\gamma$
        \item Backward propagators: $\mathcal{B}^{(\ell \to p)}_{\mu} = \prod_{k=\ell}^{p} \frac{(W^{(k+1)})^T \sigma'_{k+1}(x_\mu)}{\sqrt{m}}$
        \item Backward source terms: $\mathcal{U}^{(p)}_{\gamma\mu} = \frac{1}{m} G^{(p+1)}_\gamma \langle x^{(p)}_\gamma, x^{(p)}_\mu \rangle$
    \end{itemize}
    
The operator $\delta_\mu$ acts on any parameter-dependent quantity $F$. When $F$ is a product of terms, $\delta_\mu(F_1 F_2) = (\delta_\mu F_1)F_2 + F_1(\delta_\mu F_2)$. For ReLU networks ($\sigma''=0$), the fundamental rules are:
\begin{itemize}
    \item On weights: $\delta_\mu(W^{(k)}) = G^{(k)}_\mu (x^{(k-1)}_\mu)^T$ and $\delta_\mu((W^{(k)})^T) = x^{(k-1)}_\mu (G^{(k)}_\mu)^T$.
    \item On the output layer: $\delta_\mu(a_t) = x^{(H)}_\mu$.
    \item On activation derivatives: $\delta_\mu(\sigma'_k(x_\nu)) = 0$.
\end{itemize}

Specifically, the rewriting rules coming from \cite{huang2019dynamics} for these terms are:

\begin{itemize}
\item $a_t\rightarrow x_\mu^{(H)}$

\item $\frac{W_t^{(\ell)}}{\sqrt m}\rightarrow \diag\left(\sigma_\ell'(x_\mu) \frac{(W_t^{(\ell+1)})^\top}{\sqrt m}\cdots \sigma_H' (x_\mu)\frac{a_t}{\sqrt m}\right)\frac{\bm 1}{\sqrt m} \otimes (x_\mu^{(\ell-1)})^\top$

\item $\frac{(W_t^{(\ell)})^\top}{\sqrt m}\rightarrow  \frac{1}{\sqrt m}x_\mu^{(\ell-1)}\otimes \left(\frac{a_t^\top}{\sqrt m}\sigma_H'(x_\mu) \cdots \frac{W_t^{(\ell+1)}}{\sqrt m} \sigma_\ell'(x_\mu)\right)$

\item $x^{(\ell)}_\alpha\rightarrow \sum_{k=1}^\ell \diag\left(\sigma'_\ell(x_\alpha) \frac{W_t^{(\ell)}}{\sqrt m}\cdots \frac{W_t^{(k+1)}}{\sqrt m}\sigma_k'(x_\alpha) \sigma_k'(x_\mu)\frac{(W_t^{(k+1)})^\top}{\sqrt m}\cdots \sigma_H'(x_\mu) \frac{a_t}{\sqrt m}\right)\frac{\bm 1}{\sqrt m} \langle x_\alpha^{(k-1)}, x_\mu^{(k-1)}\rangle$

\item $\sigma_\ell^{(r)}(x_\alpha)\rightarrow\sigma^{(r+1)}_\ell(x_\alpha)\diag\left(\sigma'_\ell(x_\mu)\frac{(W_t^{(\ell+1)})^\top}{\sqrt m}\cdots \sigma_H'(x_\mu) \frac{a_t}{\sqrt m}\right) \langle x_\alpha^{(\ell-1)}, x_\mu^{(\ell-1)}\rangle + \sum_{k=1}^{\ell-1}\sigma^{(r+1)}_\ell(x_\alpha)\diag\left(\frac{W_t^{(\ell)}}{\sqrt m}\sigma_{\ell-1}'(x_\alpha)\cdots \frac{W_t^{(k+1)}}{\sqrt m}\sigma_k'(x_\alpha)\sigma'_k(x_\mu)\frac{(W_t^{(k+1)})^\top}{\sqrt m}\cdots \sigma_H'(x_\mu) \frac{a_t}{\sqrt m}\right) \langle x_\alpha^{(k-1)}, x_\mu^{(k-1)}\rangle$
\end{itemize}



From these, we derive the action on the composite objects (for any index $\nu$):
\begin{itemize}
    \item $\delta_\mu x^{(p)}_\nu = \sum_{j=1}^{p} \mathcal{J}^{(p \to j)}_{\nu} \frac{\sigma'_{j}(x_\nu)}{\sqrt{m}} \langle x^{(j-1)}_\mu, x^{(j-1)}_\nu \rangle G^{(j)}_\mu$
    \item $\delta_\mu G^{(\ell)}_\nu = \sum_{p=\ell}^{H-1} \mathcal{B}^{(\ell \to p)}_{\nu} \frac{G^{(p+1)}_\mu \langle x^{(p)}_\mu, x^{(p)}_\nu \rangle}{m} + \mathcal{B}^{(\ell \to H-1)}_{\nu} \frac{x^{(H)}_\mu}{\sqrt{m}}$
    \item Action on forward propagators $\mathcal{J}^{(p \to j)}_{\nu}$:
    \begin{align}
        \delta_\mu(\mathcal{J}^{(p \to j)}_{\nu}) = \sum_{i=j+1}^{p} \left( \prod_{k=i+1}^{p} \frac{\sigma'_{k}(x_\nu) W^{(k)}}{\sqrt{m}} \right) \frac{\sigma'_{i}(x_\nu) G^{(i)}_\mu (x^{(i-1)}_\mu)^T}{\sqrt{m}} \left( \prod_{k=j+1}^{i-1} \frac{\sigma'_{k}(x_\nu) W^{(k)}}{\sqrt{m}} \right)
    \end{align}
    \item Action on backward propagators $\mathcal{B}^{(\ell \to p)}_{\nu}$:
    \begin{align}
        \delta_\mu(\mathcal{B}^{(\ell \to p)}_{\nu}) = \sum_{i=\ell}^{p} \left( \prod_{k=i+1}^{p} \frac{(W^{(k+1)})^T \sigma'_{k+1}(x_\nu)}{\sqrt{m}} \right) \frac{x^{(i)}_\mu (G^{(i+1)}_\mu)^T \sigma'_{i+1}(x_\nu)}{\sqrt{m}} \left( \prod_{k=\ell}^{i-1} \frac{(W^{(k+1)})^T \sigma'_{k+1}(x_\nu)}{\sqrt{m}} \right)
    \end{align}
    \item Action on forward source terms $\mathcal{T}^{(j)}_{\gamma\nu}$:
    \begin{align}
        \delta_\mu(\mathcal{T}^{(j)}_{\gamma\nu}) = \frac{\sigma'_{j}(x_\nu)}{\sqrt{m}} \Biggl[ & \left( \langle \delta_\mu x^{(j-1)}_\gamma, x^{(j-1)}_\nu \rangle + \langle x^{(j-1)}_\gamma, \delta_\mu x^{(j-1)}_\nu \rangle \right) G^{(j)}_\gamma \nonumber \\
        & + \langle x^{(j-1)}_\gamma, x^{(j-1)}_\nu \rangle (\delta_\mu G^{(j)}_\gamma) \Biggr]
    \end{align}
    \item Action on backward source terms $\mathcal{U}^{(p)}_{\gamma\nu}$:
    \begin{align}
        \delta_\mu(\mathcal{U}^{(p)}_{\gamma\nu}) = \frac{1}{m} \Biggl[ & (\delta_\mu G^{(p+1)}_\gamma) \langle x^{(p)}_\gamma, x^{(p)}_\nu \rangle \nonumber \\
        & + G^{(p+1)}_\gamma \left( \langle \delta_\mu x^{(p)}_\gamma, x^{(p)}_\nu \rangle + \langle x^{(p)}_\gamma, \delta_\mu x^{(p)}_\nu \rangle \right) \Biggr]
    \end{align}
\end{itemize}

\subsection{Derivation of $K^{(4)}_{\alpha\beta\gamma\mu}$}
We apply $\delta_\mu$ to each of the three main components of $K^{(3)}_{\alpha\beta\gamma}$.

\subsubsection{Part 1: Differentiating $K^{(3, \text{out})}_{\alpha\beta\gamma}$}
The first component is $K^{(3, \text{out})}_{\alpha\beta\gamma} = \sum_{j=1}^{H} \Biggl( \left\langle \left( \prod_{k=j+1}^{H} \frac{\sigma'_{k}(x_\alpha) W^{(k)}}{\sqrt{m}} \right) \frac{\sigma'_{j}(x_\alpha)}{\sqrt{m}} \langle x^{(j-1)}_\gamma, x^{(j-1)}_\alpha \rangle G^{(j)}_\gamma, x^{(H)}_\beta \right\rangle + \left\langle x^{(H)}_\alpha, \left( \prod_{k=j+1}^{H} \frac{\sigma'_{k}(x_\beta) W^{(k)}}{\sqrt{m}} \right) \frac{\sigma'_{j}(x_\beta)}{\sqrt{m}} \langle x^{(j-1)}_\gamma, x^{(j-1)}_\beta \rangle G^{(j)}_\gamma \right\rangle \Biggr)$. Applying $\delta_\mu$ yields four terms:
\begin{align}
\delta_\mu(K^{(3, \text{out})}) = & \sum_{j=1}^{H} \Biggl( \left\langle \delta_\mu \left( \left( \prod_{k=j+1}^{H} \frac{\sigma'_{k}(x_\alpha) W^{(k)}}{\sqrt{m}} \right) \frac{\sigma'_{j}(x_\alpha)}{\sqrt{m}} \langle x^{(j-1)}_\gamma, x^{(j-1)}_\alpha \rangle G^{(j)}_\gamma \right), x^{(H)}_\beta \right\rangle \nonumber \\
& + \left\langle \left( \prod_{k=j+1}^{H} \frac{\sigma'_{k}(x_\alpha) W^{(k)}}{\sqrt{m}} \right) \frac{\sigma'_{j}(x_\alpha)}{\sqrt{m}} \langle x^{(j-1)}_\gamma, x^{(j-1)}_\alpha \rangle G^{(j)}_\gamma, \delta_\mu x^{(H)}_\beta \right\rangle \nonumber \\
& + \left\langle \delta_\mu x^{(H)}_\alpha, \left( \prod_{k=j+1}^{H} \frac{\sigma'_{k}(x_\beta) W^{(k)}}{\sqrt{m}} \right) \frac{\sigma'_{j}(x_\beta)}{\sqrt{m}} \langle x^{(j-1)}_\gamma, x^{(j-1)}_\beta \rangle G^{(j)}_\gamma \right\rangle \nonumber \\
& + \left\langle x^{(H)}_\alpha, \delta_\mu \left( \left( \prod_{k=j+1}^{H} \frac{\sigma'_{k}(x_\beta) W^{(k)}}{\sqrt{m}} \right) \frac{\sigma'_{j}(x_\beta)}{\sqrt{m}} \langle x^{(j-1)}_\gamma, x^{(j-1)}_\beta \rangle G^{(j)}_\gamma \right) \right\rangle \Biggr)
\end{align}

\subsubsection{Part 2: Differentiating $K^{(3, G)}_{\alpha\beta\gamma}$}
For each $\ell$, we differentiate:
\begin{align}
K^{(3,G)}_{\alpha\beta\gamma} = & \sum_{\ell=1}^{H} \langle x^{(\ell-1)}_\alpha, x^{(\ell-1)}_\beta \rangle \Biggl[ \nonumber \\
& \sum_{p=\ell}^{H-1} \Biggl( \left\langle \left( \prod_{k=\ell}^{p} \frac{(W^{(k+1)})^T \sigma'_{k+1}(x_\alpha)}{\sqrt{m}} \right) \frac{G^{(p+1)}_\gamma \langle x^{(p)}_\gamma, x^{(p)}_\alpha \rangle}{m}, G^{(\ell)}_\beta \right\rangle \nonumber \\
& + \left\langle G^{(\ell)}_\alpha, \left( \prod_{k=\ell}^{p} \frac{(W^{(k+1)})^T \sigma'_{k+1}(x_\beta)}{\sqrt{m}} \right) \frac{G^{(p+1)}_\gamma \langle x^{(p)}_\gamma, x^{(p)}_\beta \rangle}{m} \right\rangle \Biggr) \nonumber \\
& + \left\langle \left( \prod_{k=\ell}^{H-1} \frac{(W^{(k+1)})^T \sigma'_{k+1}(x_\alpha)}{\sqrt{m}} \right) \frac{x^{(H)}_\gamma}{\sqrt{m}}, G^{(\ell)}_\beta \right\rangle \nonumber \\
& + \left\langle G^{(\ell)}_\alpha, \left( \prod_{k=\ell}^{H-1} \frac{(W^{(k+1)})^T \sigma'_{k+1}(x_\beta)}{\sqrt{m}} \right) \frac{x^{(H)}_\gamma}{\sqrt{m}} \right\rangle \Biggr]
\end{align}

\subsubsection{Part 3: Differentiating $K^{(3, x)}_{\alpha\beta\gamma}$}
For each $\ell$, we differentiate:
\begin{align}
K^{(3,x)}_{\alpha\beta\gamma} = & \sum_{\ell=1}^{H} \langle G^{(\ell)}_\alpha, G^{(\ell)}_\beta \rangle \sum_{j=1}^{\ell-1} \Biggl( \left\langle \left( \prod_{k=j+1}^{\ell-1} \frac{\sigma'_{k}(x_\alpha) W^{(k)}}{\sqrt{m}} \right) \frac{\sigma'_{j}(x_\alpha)}{\sqrt{m}} \langle x^{(j-1)}_\gamma, x^{(j-1)}_\alpha \rangle G^{(j)}_\gamma, x^{(\ell-1)}_\beta \right\rangle \nonumber \\
& + \left\langle x^{(\ell-1)}_\alpha, \left( \prod_{k=j+1}^{\ell-1} \frac{\sigma'_{k}(x_\beta) W^{(k)}}{\sqrt{m}} \right) \frac{\sigma'_{j}(x_\beta)}{\sqrt{m}} \langle x^{(j-1)}_\gamma, x^{(j-1)}_\beta \rangle G^{(j)}_\gamma \right\rangle \Biggr)
\end{align}

Each term such as $\delta_\mu x, \delta_\gamma x, \delta_\mu G, \delta_\gamma G, \delta_\mu(\delta_\gamma x), \delta_\mu(\delta_\gamma G)$ must then be substituted by its own full non-recursive formula to get the final expression. The result is a sum of hundreds of terms, each being a product of various inner products of activations and backpropagated vectors, representing the full complexity of the fourth-order dynamics.






\newpage

\section{Full Expression for $K^{(4)}$ via Directional Derivative}

We now derive the complete, non-recursive formula for $K^{(4)}_{\alpha\beta\gamma\mu}$. We apply the directional derivative operator $\delta_\mu$ to the expression for $K^{(3)}_{\alpha\beta\gamma}$. The operator $\delta_\mu$ acts as a derivative and obeys the product rule. Since the expression for $K^{(3)}_{\alpha\beta\gamma}$ is a sum of inner products, we apply $\delta_\mu$ to each term, which expands as $\delta_\mu \langle A, B \rangle = \langle \delta_\mu A, B \rangle + \langle A, \delta_\mu B \rangle$.

The calculation involves computing $\delta_\mu$ of the core components: activations $x^{(p)}_\nu$, backpropagated vectors $G^{(\ell)}_\nu$, and the propagators which are products of weights and activation derivatives. For ReLU networks, any term involving a derivative of $\sigma'$ vanishes.

The final expression for $K^{(4)}_{\alpha\beta\gamma\mu} = \delta_\mu(K^{(3)}_{\alpha\beta\gamma})$ is the sum of the following four components, corresponding to the application of $\delta_\mu$ to each of the four main terms in the $K^{(3)}$ formula.

\begin{enumerate}
    \item \textbf{Term 1: $\delta_\mu$ on the first part of $K^{(3, \text{out})}$ -- $\langle \delta_\gamma x^{(H)}_\alpha, x^{(H)}_\beta \rangle$}
    
    This term expands to $\langle \delta_\mu(\delta_\gamma x^{(H)}_\alpha)), x^{(H)}_\beta \rangle + \langle \delta_\gamma x^{(H)}_\alpha, \delta_\mu x^{(H)}_\beta \rangle$. The second part, $\langle \delta_\gamma x^{(H)}_\alpha, \delta_\mu x^{(H)}_\beta \rangle$, is a symmetric product of first-order variations. The first part, involving $\delta_\mu(\delta_\gamma x^{(H)}_\alpha))$, gives rise to second-order variations. This yields:
    \begin{align}
    & \sum_{j=1}^{H} \Biggl( \left\langle \delta_\mu \left( \mathcal{J}^{(H \to j)}_{\alpha} \mathcal{T}^{(j)}_{\gamma\alpha} \right), x^{(H)}_\beta \right\rangle + \left\langle \mathcal{J}^{(H \to j)}_{\alpha} \mathcal{T}^{(j)}_{\gamma\alpha}, \delta_\mu x^{(H)}_\beta \right\rangle \Biggr)
    \end{align}
    where $\mathcal{J}^{(H \to j)}_{\alpha} \mathcal{T}^{(j)}_{\gamma\alpha} = \delta_\gamma x^{(H)}_\alpha$.

    \item \textbf{Term 2: $\delta_\mu$ on the second part of $K^{(3, \text{out})}$ -- $\langle x^{(H)}_\alpha, \delta_\gamma x^{(H)}_\beta \rangle$}
    
    Symmetric to the first term, this expands to $\langle \delta_\mu x^{(H)}_\alpha, \delta_\gamma x^{(H)}_\beta \rangle + \langle x^{(H)}_\alpha, \delta_\mu(\delta_\gamma x^{(H)}_\beta) \rangle$.
    \begin{align}
    & \sum_{j=1}^{H} \Biggl( \left\langle \delta_\mu x^{(H)}_\alpha, \mathcal{J}^{(H \to j)}_{\beta} \mathcal{T}^{(j)}_{\gamma\beta} \right\rangle + \left\langle x^{(H)}_\alpha, \delta_\mu \left( \mathcal{J}^{(H \to j)}_{\beta} \mathcal{T}^{(j)}_{\gamma\beta} \right) \right\rangle \Biggr)
    \end{align}

    \item \textbf{Term 3: $\delta_\mu$ on the $K^{(3, G)}$ terms}

    Here, $\delta_\mu$ acts on the term structure $\langle \dots \rangle \left( \langle \delta_\gamma G, G \rangle + \langle G, \delta_\gamma G \rangle \right)$. This involves differentiating the outer product of activations $\langle x^{(\ell-1)}_\alpha, x^{(\ell-1)}_\beta \rangle$ and the inner products of backpropagated vectors.
    \begin{align}
    & \sum_{\ell=1}^{H} \left( \langle \delta_\mu x^{(\ell-1)}_\alpha, x^{(\ell-1)}_\beta \rangle + \langle x^{(\ell-1)}_\alpha, \delta_\mu x^{(\ell-1)}_\beta \rangle \right) \left[ \dots \right] \nonumber \\
    & + \sum_{\ell=1}^{H} \langle x^{(\ell-1)}_\alpha, x^{(\ell-1)}_\beta \rangle \left[ \delta_\mu(\dots) \right]
    \end{align}
    where the bracket contains the full expression from $K^{(3,G)}$. For example, a term $\langle \delta_\gamma G^{(\ell)}_\alpha, G^{(\ell)}_\beta \rangle$ becomes $\langle \delta_\mu(\delta_\gamma G^{(\ell)}_\alpha), G^{(\ell)}_\beta \rangle + \langle \delta_\gamma G^{(\ell)}_\alpha, \delta_\mu G^{(\ell)}_\beta \rangle$.

    \item \textbf{Term 4: $\delta_\mu$ on the $K^{(3, x)}$ terms}

    Similarly, $\delta_\mu$ acts on $\langle G, G \rangle \left( \langle \delta_\gamma x, x \rangle + \langle x, \delta_\gamma x \rangle \right)$, differentiating the $\langle G,G \rangle$ factor and the inner factor involving activations.
    \begin{align}
    & \sum_{\ell=1}^{H} \left( \langle \delta_\mu G^{(\ell)}_\alpha, G^{(\ell)}_\beta \rangle + \langle G^{(\ell)}_\alpha, \delta_\mu G^{(\ell)}_\beta \rangle \right) \left[ \dots \right] \nonumber \\
    & + \sum_{\ell=1}^{H} \langle G^{(\ell)}_\alpha, G^{(\ell)}_\beta \rangle \left[ \delta_\mu(\dots) \right]
    \end{align}
\end{enumerate}

Expanding these terms fully using the rewriting rules for $\delta_\mu$ on each constituent component ($W, a_t, x, G$) leads to an exceedingly large but well-defined expression, representing the complete fourth-order kernel. The structure involves symmetric terms and nested applications of the derivative operators, with each $\delta$ operator adding a sum over the layers of the network. The full expansion is omitted for brevity but its construction is a direct application of the rules outlined.

\section{Recursive Formula for $K^{(3)}$ as a Function of Depth $H$}

To analyze the influence of network depth, it is useful to establish a recursive relation for $K^{(3)}$. Let $K^{(3, H)}_{\alpha\beta\gamma}$ denote the kernel for a network with $H$ layers. The goal is to express $K^{(3, H+1)}_{\alpha\beta\gamma}$ in terms of quantities computed for a network with $H$ layers.

The starting point is the decomposition of $K^{(2, H+1)}_{\alpha\beta}$, the NTK for a network with $H+1$ layers:
\begin{equation}
K^{(2, H+1)}_{\alpha\beta} = \underbrace{\left( \langle x^{(H)}_\alpha, x^{(H)}_\beta \rangle + \sum_{\ell=1}^{H} \langle G'^{(\ell)}_\alpha, G'^{(\ell)}_\beta \rangle \langle x^{(\ell-1)}_\alpha, x^{(\ell-1)}_\beta \rangle \right)}_{K^{(2,H)}_{\alpha\beta} \text{ with } a'_t} + T_{H+1} + C_{H+1}
\end{equation}
where:
\begin{itemize}
    \item $x^{(H+1)}_\mu = \frac{1}{\sqrt{m}}\sigma(W^{(H+1)} x^{(H)}_\mu)$.
    \item $a'_t = \frac{1}{\sqrt{m}} (W^{(H+1)})^T \sigma'_{H+1} a_t$ is an "effective output vector" that replaces $a_t$ in the computations of $G^{(\ell)}$ for $\ell \le H$.
    \item $T_{H+1} = \langle G^{(H+1)}_\alpha, G^{(H+1)}_\beta \rangle \langle x^{(H)}_\alpha, x^{(H)}_\beta \rangle$ is the term from the new layer $H+1$.
    \item $C_{H+1} = \langle x^{(H+1)}_\alpha, x^{(H+1)}_\beta \rangle - \langle a'_t, a'_t \rangle \langle x^{(H)}_\alpha, x^{(H)}_\beta \rangle$ is a correction term.
\end{itemize}
The vectors $G'^{(\ell)}$ are the backpropagated gradients in a network with $H$ layers but with $a'_t$ as output weights, while $G^{(H+1)}$ is the gradient for layer $H+1$.

Applying the directional derivative operator $\delta_\gamma(\cdot) = \langle \nabla_\theta (\cdot), \nabla_\theta f_t(x_\gamma) \rangle$, we obtain the recursion for $K^{(3)}$:
\begin{equation}
K^{(3, H+1)}_{\alpha\beta\gamma} = \delta_\gamma \left(K^{(2,H)}_{\alpha\beta, a'_t} \right) + \delta_\gamma(T_{H+1}) + \delta_\gamma(C_{H+1})
\end{equation}
The term $\delta_\gamma \left(K^{(2,H)}_{\alpha\beta, a'_t} \right)$ is analogous to $K^{(3,H)}_{\alpha\beta\gamma}$, but with additional dependencies through $a'_t$. The other terms expand as follows:
\begin{align}
\delta_\gamma(T_{H+1}) &= \langle \delta_\gamma G^{(H+1)}_\alpha, G^{(H+1)}_\beta \rangle \langle x^{(H)}_\alpha, x^{(H)}_\beta \rangle + \dots \\
\delta_\gamma(C_{H+1}) &= \langle \delta_\gamma x^{(H+1)}_\alpha, x^{(H+1)}_\beta \rangle - \langle \delta_\gamma a'_t, a'_t \rangle \langle x^{(H)}_\alpha, x^{(H)}_\beta \rangle + \dots
\end{align}
This decomposition allows us to separate the effect of adding a layer by isolating terms specific to the new layer $(T_{H+1}, C_{H+1})$ and the modification of previous layer terms through $a'_t$.

\newpage


\bibliographystyle{plain}
\bibliography{references}

\appendix
\section{References}

\bibitem{dyer2019asymptoticswidenetworksfeynman}
E. Dyer and G. Gur-Ari,
``Asymptotics of Wide Networks from Feynman Diagrams,''
arXiv:1909.11304 [cs.LG], 2019.

\bibitem{huang2019dynamics}
J. Huang and H.-T. Yau,
``Dynamics of Deep Neural Networks and Neural Tangent Hierarchy,''
arXiv:1909.08156 [cs.LG], 2019.

\bibitem{terjek2025mlpseocconcentrationntk}
D. Terjék and D. González-Sánchez,
``MLPs at the EOC: Concentration of the NTK,''
arXiv:2501.14724 [cs.LG], 2025.

\end{document}